\scanpage{64}
\mspage{II. 12}
% lines 1-2, cancelled
\deleted{\unsure{constant tordu} sur $X$, dont la formation commute à \unsure{tout changement de la partie connexe à zéro}}

% lines 3-10
et que l'énoncé analogue est vrai après réduction de la situation à un corps de base de type fini. Cette hypothèse est \inserted{donc} inutile si on suppose qu'on dispose de la résolution des singularités pour les schémas de type fini \inserted{de dim $\leq\dim X_{\unsure{\overline\eta|\overline y}}$,} sur une clôture algébrique de $K$ ($\eta$ étant le pt générique de $Y$), \deleted{donc} \deleted{\unsure{elle est}} comparez I 8.8, \unsure{8.9 b}. Il en est en particulier ainsi si $\operatorname{car} k=0$, ou si $\dim X_\eta\leq2$.

% lines 11-16
On peut montrer que la conclusion est également valable si $i=1$, sans restriction \deleted{sur} de propreté sur $f$, \deleted{\unsure{dans ce cas cf}} en utilisant les réductions faites dans les n\textsuperscript{os} suivants; comparez n\textsuperscript{o} 5.

% lines 17-31, the whole remark crossed by diagonal strokes
\deleted{\uline{Remarque} 1.17. \deleted{Il est plausible que la forme} Avec les notations de 1.1, \deleted{\unsure{étant}} \deleted{\unsure{généralisée}} et soit $H$ un groupe \inserted{profini} quotient de $I$, invariant par l'action de $\pi$ par automorphismes intérieurs, et tel que $H$ soit un \deleted{\unsure{groupe analytique}} \deleted{\unsure{$\mathbb{Q}_\ell$}} \deleted{$\ell$-adique résoluble} pro-$\ell$-groupe \uuline{résoluble}. On trouve donc encore une extension}
\[
\deleted{1\mathrel{\text{---}}H\mathrel{\text{---}}\pi'\longrightarrow\pi_0\mathrel{\text{---}}1}.
\]
\deleted{Ceci posé, \uline{il est} plausible que l'on ait encore $\overline{[\pi_0,H]}$ ouvert dans $\pi'$. Cela semblerait devoir résulter de l'hyp. de Riemann-Weil en dim. quelconque, comme 1.1. résulte de l'hyp. de R.W. en dim 1. S'il en était ainsi, on}
% archive sheet number: 63
